% =========================================================
% sec_axioms_A.tex — A/Аксиоматика (вставка) — АКТУАЛЬНАЯ РЕДАКЦИЯ
% Назначение: строгие определения/аксиомы/постулаты, на которые
% ссылаются разделы математики (B) и верификации (C).
% Подключение: % =========================================================
% sec_axioms_A.tex — A/Аксиоматика (вставка) — АКТУАЛЬНАЯ РЕДАКЦИЯ
% Назначение: строгие определения/аксиомы/постулаты, на которые
% ссылаются разделы математики (B) и верификации (C).
% Подключение: % =========================================================
% sec_axioms_A.tex — A/Аксиоматика (вставка) — АКТУАЛЬНАЯ РЕДАКЦИЯ
% Назначение: строгие определения/аксиомы/постулаты, на которые
% ссылаются разделы математики (B) и верификации (C).
% Подключение: % =========================================================
% sec_axioms_A.tex — A/Аксиоматика (вставка) — АКТУАЛЬНАЯ РЕДАКЦИЯ
% Назначение: строгие определения/аксиомы/постулаты, на которые
% ссылаются разделы математики (B) и верификации (C).
% Подключение: \input{sec_axioms_A}
% Совместимость: article, pdfLaTeX, T2A, babel (ru/uk)
% =========================================================

\section{Аксиоматика и определения A (расширение канона)}

\subsection*{Статус}
\textbf{Ответственный:} Сотрудник A (аксиоматика) \quad
\textbf{Статус:} рабочая вставка \quad
\textbf{Версия:} A-0.2 \quad
\textbf{Дата (Europe/Kyiv):} \today

% ---------------------------------------------------------
\subsection{A.1. Термины и фазовые состояния}

\begin{definition}[Фазовые состояния эфира]
В модели различаются два фазовых состояния движения Амеров:
\begin{itemize}
\item \textbf{Хаотическая фаза:} статистически хаотическое движение Амеров.
Вне \EVO\ хаотическая фаза называется \SE\ (Свободный эфир), а внутри \EVO\ — \NE\ (Несвободный эфир),
включая центральное ядро \NY\ и межканальные прослойки.
\item \textbf{Направленная фаза:} упорядоченное циркуляционное движение, существующее \emph{только} в каналах \EVO.
Направленное движение в каналах \textbf{не} называется \NE.
\end{itemize}
\end{definition}

\begin{axiom}[Запрет устойчивых потоков \SE]
В \SE\ не допускаются устойчивые потоки как первичная онтология.
Допускаются только распространяющиеся \VR\ как флуктуации импульса/плотности,
передаваемые цепочками парных соударений.
\end{axiom}

% ---------------------------------------------------------
\subsection{A.2. Геометрия \EVO: каналы и граничные поверхности}

\begin{definition}[Каналы и граничные поверхности (\GP)]
\EVO\ состоит из одного или нескольких коаксиальных циркуляционных каналов (направленная фаза)
и областей \NE\ (хаос) между ними и в \NY.
Каждый канал имеет две \GP: внутреннюю и внешнюю.
Внешняя \GP внешнего канала является внешней \GP всего \EVO\ (границей \EVO\ с \SE).
\end{definition}

\begin{remark}[Шнековая организация \GP]
Шнекоподобная форма \GP понимается как согласованная пространственно-временная организация контактов,
которая может быть бегущей или стоячей (см. A.6) и влияет на статистику соударений на \GP.
\end{remark}

% ---------------------------------------------------------
\subsection{A.3. Инвариант и жёсткие запреты}

\begin{axiom}[Единственный физический акт]
Единственный физический акт модели — парные упругие соударения Амеров.
Все процессы (структурирование, перенос \VR, взаимодействия, движение/вращение \EVO) реализуются только через них.
\end{axiom}

\begin{axiom}[Жёсткие запреты модели]
Запрещается вводить как первичные сущности:
\begin{itemize}
\item устойчивые потоки \SE;
\item ``давление \SE'' как первичный агент;
\item проникновение Амеров через \GP;
\item поглощение/переизлучение Амеров;
\item ``виртуальные переносчики'' как физические объекты;
\item ``поле'' как отдельную субстанцию (термин допустим только как макро-ярлык).
\end{itemize}
Генезис вихрей из \SE\ в рамках канона не рассматривается (\EVO\ считаются существующими).
\end{axiom}

% ---------------------------------------------------------
\subsection{A.4. Условие целостности хаотической фазы (запрет на обгон)}

\begin{postulate}[Целостность хаотической фазы у \GP]
В хаотической фазе (\SE\ и \NE) у \GP поддерживается режим, обеспечивающий целостность \EVO:
\[
l_{\max}\le d_A, \qquad \mean{l}=0.5\,d_A,
\]
где \(l\) — длина свободного пробега Амера, \(d_A\) — диаметр Амера.
При нарушении режима целостности вихревая структура теряет устойчивость и распадается.
\end{postulate}

\begin{remark}[Каналы]
В каналах (направленная фаза) концентрация ниже, скорости выше;
обгоны/отставания возможны внутри потока, но отсутствует движение против общего направления.
\end{remark}

% ---------------------------------------------------------
\subsection{A.5. КУБК: контактный баланс и минимальная концентрация канала}

\begin{postulate}[КУБК: баланс \emph{числа} контактов на \GP]
В квазистационарном режиме на каждом участке \GP отсутствует устойчивый перекос по \emph{числу} контактов
между амерами хаотической фазы (\SE/\NE) и амерами канала у \GP.
Эквивалентно, после осреднения по времени выполняется согласование \emph{интенсивностей контактов (касания \GP)}:
\[
\Phi^-_{\mathrm{хаос}}(x)=\Phi^+_{\mathrm{канал}}(x),\qquad x\in\Sigma,
\]
где \(\Sigma\) — соответствующая \GP, \(n\) — внешняя единичная нормаль к \(\Sigma\),
а \(\Phi\) не является ``потоком \SE\ в объёме'', а есть локальная величина \emph{частоты/интенсивности касаний поверхности}:
\[
\Phi^-_{\mathrm{хаос}}(x)=\int_{v\cdot n<0}\abs{v\cdot n}\;f_{\mathrm{хаос}}(x,v)\,d^3v,\qquad
\Phi^+_{\mathrm{канал}}(x)=\int_{v\cdot n>0}(v\cdot n)\;f_{\mathrm{канал}}(x,v)\,d^3v.
\]
Здесь \(f_{\mathrm{хаос}}(x,v)\) и \(f_{\mathrm{канал}}(x,v)\) — локальные функции распределения амеров по скоростям
в окрестности \(\Sigma\) со своей стороны поверхности; интегрирование ведётся по соответствующим полупространствам
\(v\cdot n<0\) и \(v\cdot n>0\).
\end{postulate}

\begin{postulate}[Принцип минимальной концентрации канала у \GP]
Среди всех режимов канала, совместимых с целостностью \EVO, реализуется режим с \emph{минимально возможной}
локальной концентрацией \(n_{\mathrm{канал}}(x)\) у \GP, при котором сохраняется КУБК.
\end{postulate}

\begin{remark}[Различение балансов]
КУБК фиксирует баланс \emph{числа} контактов.
Дрейф/вращение определяются не этим балансом, а анизотропией \emph{приращений импульса} в контактах на \GP.
\end{remark}

% ---------------------------------------------------------
\subsection{A.6. Динамическая природа ``шнека''}

\begin{postulate}[Шнек как колебательный режим \GP]
Равновесие на \GP поддерживается колебательным процессом границы относительно среднего положения.
Колебания \GP локально ортогональны направлению основного потока в канале.
Период/фаза колебаний связаны (или фазно согласованы) с характерным временем циркуляции в структуре канала,
в частности с временем обращения вокруг малого кольца тора и/или его гармониками.
Возможны режимы: стоячий шнек и бегущий шнек (по ходу или против хода канального движения),
в зависимости от распределения внутренних/межканальных \VR\ и состояния \KP.
\end{postulate}

% ---------------------------------------------------------
\subsection{A.7. Симметрия дрейфа, хиральность и двуторий}

\begin{postulate}[Симметрия дрейфа в изотропном фоне]
В изотропном фоне (\(\nabla\KP=0\)) и при отсутствии внешних \VR, выделяющих направление,
идеализированный симметричный ``тонкий'' тор не имеет выделенного направления для устойчивого дрейфа.
Следовательно, дрейф требует нарушения симметрии статистики \emph{приращений импульса} на \GP
(внешнего или самосогласованного внутреннего).
\end{postulate}

\begin{definition}[Хиральность (зарученность)]
Хиральность \(\chi\) — дискретный паспортный знак зарученности циркуляции канала/объекта:
\[
\chi\in\{-1,+1\}.
\]
\(\chi\) используется как метка конфигурации (паспорт вихря) и не вводится как отдельная физическая субстанция.
\end{definition}

\begin{definition}[Двуторий]
Двуторий — тороидально-спиральная форма \EVO, возникающая как самосогласованная трансформация ``тонкого'' вихря,
когда собственные \VR\ и связанная с ними структура \KP\ внутри \EVO\ приводят к устойчивому нарушению симметрии
статистики контактов/приращений импульса на \GP.
\end{definition}

\begin{postulate}[Двуторий: механизм ``вращение \(\to\) дрейф'']
В двутории внутренняя топология и шнековая организация \GP обеспечивают преобразование вращательных мод
в поступательный дрейф вдоль оси симметрии.
Механизм трактуется кинетически: интерференция собственных \VR\ формирует устойчивую анизотропию состояния хаотической фазы
(градиенты \KP) у разных участков \GP; интеграл микросмещений мест соударений на \GP даёт ненулевой дрейф и/или момент.
\end{postulate}

\begin{remark}
Формализация дрейфа/момента через поверхностные интегралы от анизотропии импульсного обмена \(j(x,t)\)
является предметом раздела B/Математика.
\end{remark}

% ---------------------------------------------------------
\subsection{A.8. Прокси (K, G, A\_ij): измерительные ярлыки (норматив)}

\begin{axiom}[Прокси как язык измерений, не сущности]
В целях сопоставимости расчётов, симуляций и эксперимента допускается использование прокси-наборов
\((K,\,G,\,A_{ij})\) как \emph{измерительных ярлыков} для состояния хаотической фазы в окрестности \(\Sigma\)
(т.е. для канонических понятий \(\KP\) и \(\VR\)).

\textbf{Нормативно:} прокси не являются новыми сущностями и не добавляют причинности в модель.
Причинность остаётся кинетической: соударения, КУБК (баланс числа контактов) и анизотропия приращений импульса на \(\Sigma\).
\end{axiom}

\begin{remark}[Запретная лексика для прокси]
Запрещены формулировки причинности вида: ``\(G\) действует/толкает'',
``тензор \(A_{ij}\) создаёт силу'', ``поле \(K\) ускоряет'' и т.п.
Допустимо только: ``прокси \(G\)/\(A_{ij}\) наблюдается/оценён'' и ``используется как измерительная метка''.
\end{remark}

\begin{remark}[Рекомендуемая нормировка (не обязательна)]
Допускается использовать единый вид для первых моментов распределения:
\[
K=n\langle v^2\rangle,\qquad G=n\langle v\rangle,\qquad
A_{ij}=P_{ij}-\frac{1}{3}\tr(P)\,\delta_{ij},\qquad P_{ij}=n\langle v_i v_j\rangle.
\]
\end{remark}

% ---------------------------------------------------------
\subsection{A.9. Нормы тестов шнека: ось e\_\parallel и разведение ``V=0''}

\begin{postulate}[Канонический выбор оси \(e_{\parallel}\)]
Для тестов шнека вводится ось объекта \(e_{\parallel}\) — \textbf{ось симметрии} тора/двутория,
вдоль которой определяется продольный дрейф.

\textbf{Нормативно:} \(e_{\parallel}\) должен быть однозначно воспроизводим из геометрии объекта и его паспорта,
без апелляции к внешним ``полям''.
Для тора/двутория ориентация \(e_{\parallel}\) согласуется с принятой параметризацией \(X(\varphi,\theta)\) поверхности и паспортом вихря:
касательное направление \(\partial_{\varphi}X\) соответствует циркуляции по большой окружности, а \(e_{\parallel}\) — оси симметрии выбранной системы координат.
\end{postulate}

\begin{postulate}[Разведение норм ``\(V=0\)'' и нормы по \(V_{\parallel}\)]
В тестах используются две разные нормы ``нуля'', их нельзя смешивать:
\begin{enumerate}
\item \textbf{Охранная норма (T2/T6):} запись ``\(V=0\)'' означает \textbf{нулевой вектор} скорости дрейфа,
т.е. ноль всех компонент:
\[
V=0 \iff V_x=V_y=V_z=0.
\]
\item \textbf{Норма шнека (T5--T6):} нормативной величиной является \textbf{продольная компонента}
\[
V_{\parallel}\equiv V\cdot e_{\parallel}.
\]
При прочих равных \(\chi\)-нечётной обязана быть именно \(V_{\parallel}\):
\[
V_{\parallel}(\chi)=-V_{\parallel}(-\chi).
\]
При фиксированном \(\chi\) реверс направления бега шнека обязан реверсировать знак \(V_{\parallel}\).
Для стоячего шнека нормативно \(V_{\parallel}=0\) (при отсутствии внешнего нарушителя).
\end{enumerate}
\end{postulate}

\begin{remark}
Смысл нормы: \(\chi\) должна входить в замыкания так, чтобы \(\chi\)-нечётная часть отклика давала вклад в \(V_{\parallel}\),
а не подменялась ``скрытым потоком \SE'' или асимметрией числа контактов.
\end{remark}

\section{Аксиоматическая рамка текущего шага: локальный патч как временное сужение задачи}
\label{sec:a_current_step}

\subsection{Глобальная цель и смысл временного сужения}

Цель текущего шага не состоит в замене глобальной тороидальной или двуториальной
геометрии ЭВО локальным цилиндрическим описанием.
Глобальная цель остаётся прежней:
выявить кинетический механизм дрейфа и ротации ЭВО через контактную статистику
и анизотропию приращений импульса на граничной поверхности $\Sigma$.

Квазицилиндрическое приближение вводится только как временное локальное сужение задачи.
Оно используется там, где режим тонкого тора $r \ll R$ позволяет считать короткий участок
большой окружности локально цилиндрическим.
Тем самым локальный патч $\Sigma_{\mathrm{loc}} \subset \Sigma$
служит не новой онтологией, а увеличительным языком записи,
удобным для выделения локальных механизмов на ГП.

\subsection{Что именно сохраняется неизменным}

В аксиоматическом смысле сохраняются следующие положения:

\begin{enumerate}
\item ЭВО глобально остаётся тороидальной или двуториальной структурой.
\item КУБК относится только к балансу числа контактов на $\Sigma$.
\item Дрейф и вращение допускаются только через анизотропию $\Delta p$ на $\Sigma$.
\item Шнек остаётся динамическим режимом ГП, а не отдельной сущностью и не потоком СЭ.
\item Локальный патч не заменяет глобальный интеграл по $\Sigma$, а только помогает его разложить.
\end{enumerate}

Следовательно, переход к $\Sigma_{\mathrm{loc}}$ не является сменой физической картины,
а является временным усилением разрешающей способности математического языка.

\subsection{Почему непрерывная q-схема оказалась переходной}

На предыдущем этапе непрерывный параметр $q$ использовался как язык локальной намотки.
Однако такой язык оказался недостаточным для сохранения физики шнекового процесса,
поскольку не различал явно:

\begin{itemize}
\item траектории Амеров в канале;
\item шнековую моду граничной поверхности;
\item дискретную связность шнековой организации с канальной навивкой.
\end{itemize}

Поэтому в текущем шаге физический приоритет непрерывного $q$ снимается.
Если параметр $q$ и сохраняется, то только как вспомогательный эффективный
локальный параметр записи $q_{\mathrm{eff}}$.

\subsection{Новая кандидатная рамка текущего шага}

В рабочем контуре проекта принимается кандидатная схема:

\[
m_A \in \mathbb{Z}_{>0},
\qquad
m_S \in \{1,2,4\},
\]

где $m_A$ задаёт квантовое число канальной навивки,
а $m_S$ задаёт базовую пространственную моду шнека на ГП.

Базовыми режимами считаются:
\[
S1,\quad S2,\quad S4,
\]
а все прочие режимы допускаются только как производные от них.

В локальном патч-языке должны различаться два касательных направления:
\[
h_A \quad \text{и} \quad h_S,
\]
где $h_A$ есть директор траектории Амеров,
а $h_S$ --- директор шнековой моды ГП.
Шнек не отождествляется с направлением канальной траектории.

\subsection{Интерпретационная граница}

Текущий шаг проекта следует понимать как последовательность:

\[
\text{глобальная цель}
\;\to\;
\text{тонкий тор}
\;\to\;
\Sigma_{\mathrm{loc}}
\;\to\;
(h_A,h_S)
\;\to\;
(m_A,m_S,S1/S2/S4)
\;\to\;
\text{возврат к глобальному интегралу по } \Sigma.
\]

Тем самым локальный акцент на квазицилиндрическом патче не замещает канон,
а служит промежуточным этапом прояснения механизма.
Именно в таком виде данный шаг должен отражаться в витрине и в журнале изменений.

% Совместимость: article, pdfLaTeX, T2A, babel (ru/uk)
% =========================================================

\section{Аксиоматика и определения A (расширение канона)}

\subsection*{Статус}
\textbf{Ответственный:} Сотрудник A (аксиоматика) \quad
\textbf{Статус:} рабочая вставка \quad
\textbf{Версия:} A-0.2 \quad
\textbf{Дата (Europe/Kyiv):} \today

% ---------------------------------------------------------
\subsection{A.1. Термины и фазовые состояния}

\begin{definition}[Фазовые состояния эфира]
В модели различаются два фазовых состояния движения Амеров:
\begin{itemize}
\item \textbf{Хаотическая фаза:} статистически хаотическое движение Амеров.
Вне \EVO\ хаотическая фаза называется \SE\ (Свободный эфир), а внутри \EVO\ — \NE\ (Несвободный эфир),
включая центральное ядро \NY\ и межканальные прослойки.
\item \textbf{Направленная фаза:} упорядоченное циркуляционное движение, существующее \emph{только} в каналах \EVO.
Направленное движение в каналах \textbf{не} называется \NE.
\end{itemize}
\end{definition}

\begin{axiom}[Запрет устойчивых потоков \SE]
В \SE\ не допускаются устойчивые потоки как первичная онтология.
Допускаются только распространяющиеся \VR\ как флуктуации импульса/плотности,
передаваемые цепочками парных соударений.
\end{axiom}

% ---------------------------------------------------------
\subsection{A.2. Геометрия \EVO: каналы и граничные поверхности}

\begin{definition}[Каналы и граничные поверхности (\GP)]
\EVO\ состоит из одного или нескольких коаксиальных циркуляционных каналов (направленная фаза)
и областей \NE\ (хаос) между ними и в \NY.
Каждый канал имеет две \GP: внутреннюю и внешнюю.
Внешняя \GP внешнего канала является внешней \GP всего \EVO\ (границей \EVO\ с \SE).
\end{definition}

\begin{remark}[Шнековая организация \GP]
Шнекоподобная форма \GP понимается как согласованная пространственно-временная организация контактов,
которая может быть бегущей или стоячей (см. A.6) и влияет на статистику соударений на \GP.
\end{remark}

% ---------------------------------------------------------
\subsection{A.3. Инвариант и жёсткие запреты}

\begin{axiom}[Единственный физический акт]
Единственный физический акт модели — парные упругие соударения Амеров.
Все процессы (структурирование, перенос \VR, взаимодействия, движение/вращение \EVO) реализуются только через них.
\end{axiom}

\begin{axiom}[Жёсткие запреты модели]
Запрещается вводить как первичные сущности:
\begin{itemize}
\item устойчивые потоки \SE;
\item ``давление \SE'' как первичный агент;
\item проникновение Амеров через \GP;
\item поглощение/переизлучение Амеров;
\item ``виртуальные переносчики'' как физические объекты;
\item ``поле'' как отдельную субстанцию (термин допустим только как макро-ярлык).
\end{itemize}
Генезис вихрей из \SE\ в рамках канона не рассматривается (\EVO\ считаются существующими).
\end{axiom}

% ---------------------------------------------------------
\subsection{A.4. Условие целостности хаотической фазы (запрет на обгон)}

\begin{postulate}[Целостность хаотической фазы у \GP]
В хаотической фазе (\SE\ и \NE) у \GP поддерживается режим, обеспечивающий целостность \EVO:
\[
l_{\max}\le d_A, \qquad \mean{l}=0.5\,d_A,
\]
где \(l\) — длина свободного пробега Амера, \(d_A\) — диаметр Амера.
При нарушении режима целостности вихревая структура теряет устойчивость и распадается.
\end{postulate}

\begin{remark}[Каналы]
В каналах (направленная фаза) концентрация ниже, скорости выше;
обгоны/отставания возможны внутри потока, но отсутствует движение против общего направления.
\end{remark}

% ---------------------------------------------------------
\subsection{A.5. КУБК: контактный баланс и минимальная концентрация канала}

\begin{postulate}[КУБК: баланс \emph{числа} контактов на \GP]
В квазистационарном режиме на каждом участке \GP отсутствует устойчивый перекос по \emph{числу} контактов
между амерами хаотической фазы (\SE/\NE) и амерами канала у \GP.
Эквивалентно, после осреднения по времени выполняется согласование \emph{интенсивностей контактов (касания \GP)}:
\[
\Phi^-_{\mathrm{хаос}}(x)=\Phi^+_{\mathrm{канал}}(x),\qquad x\in\Sigma,
\]
где \(\Sigma\) — соответствующая \GP, \(n\) — внешняя единичная нормаль к \(\Sigma\),
а \(\Phi\) не является ``потоком \SE\ в объёме'', а есть локальная величина \emph{частоты/интенсивности касаний поверхности}:
\[
\Phi^-_{\mathrm{хаос}}(x)=\int_{v\cdot n<0}\abs{v\cdot n}\;f_{\mathrm{хаос}}(x,v)\,d^3v,\qquad
\Phi^+_{\mathrm{канал}}(x)=\int_{v\cdot n>0}(v\cdot n)\;f_{\mathrm{канал}}(x,v)\,d^3v.
\]
Здесь \(f_{\mathrm{хаос}}(x,v)\) и \(f_{\mathrm{канал}}(x,v)\) — локальные функции распределения амеров по скоростям
в окрестности \(\Sigma\) со своей стороны поверхности; интегрирование ведётся по соответствующим полупространствам
\(v\cdot n<0\) и \(v\cdot n>0\).
\end{postulate}

\begin{postulate}[Принцип минимальной концентрации канала у \GP]
Среди всех режимов канала, совместимых с целостностью \EVO, реализуется режим с \emph{минимально возможной}
локальной концентрацией \(n_{\mathrm{канал}}(x)\) у \GP, при котором сохраняется КУБК.
\end{postulate}

\begin{remark}[Различение балансов]
КУБК фиксирует баланс \emph{числа} контактов.
Дрейф/вращение определяются не этим балансом, а анизотропией \emph{приращений импульса} в контактах на \GP.
\end{remark}

% ---------------------------------------------------------
\subsection{A.6. Динамическая природа ``шнека''}

\begin{postulate}[Шнек как колебательный режим \GP]
Равновесие на \GP поддерживается колебательным процессом границы относительно среднего положения.
Колебания \GP локально ортогональны направлению основного потока в канале.
Период/фаза колебаний связаны (или фазно согласованы) с характерным временем циркуляции в структуре канала,
в частности с временем обращения вокруг малого кольца тора и/или его гармониками.
Возможны режимы: стоячий шнек и бегущий шнек (по ходу или против хода канального движения),
в зависимости от распределения внутренних/межканальных \VR\ и состояния \KP.
\end{postulate}

% ---------------------------------------------------------
\subsection{A.7. Симметрия дрейфа, хиральность и двуторий}

\begin{postulate}[Симметрия дрейфа в изотропном фоне]
В изотропном фоне (\(\nabla\KP=0\)) и при отсутствии внешних \VR, выделяющих направление,
идеализированный симметричный ``тонкий'' тор не имеет выделенного направления для устойчивого дрейфа.
Следовательно, дрейф требует нарушения симметрии статистики \emph{приращений импульса} на \GP
(внешнего или самосогласованного внутреннего).
\end{postulate}

\begin{definition}[Хиральность (зарученность)]
Хиральность \(\chi\) — дискретный паспортный знак зарученности циркуляции канала/объекта:
\[
\chi\in\{-1,+1\}.
\]
\(\chi\) используется как метка конфигурации (паспорт вихря) и не вводится как отдельная физическая субстанция.
\end{definition}

\begin{definition}[Двуторий]
Двуторий — тороидально-спиральная форма \EVO, возникающая как самосогласованная трансформация ``тонкого'' вихря,
когда собственные \VR\ и связанная с ними структура \KP\ внутри \EVO\ приводят к устойчивому нарушению симметрии
статистики контактов/приращений импульса на \GP.
\end{definition}

\begin{postulate}[Двуторий: механизм ``вращение \(\to\) дрейф'']
В двутории внутренняя топология и шнековая организация \GP обеспечивают преобразование вращательных мод
в поступательный дрейф вдоль оси симметрии.
Механизм трактуется кинетически: интерференция собственных \VR\ формирует устойчивую анизотропию состояния хаотической фазы
(градиенты \KP) у разных участков \GP; интеграл микросмещений мест соударений на \GP даёт ненулевой дрейф и/или момент.
\end{postulate}

\begin{remark}
Формализация дрейфа/момента через поверхностные интегралы от анизотропии импульсного обмена \(j(x,t)\)
является предметом раздела B/Математика.
\end{remark}

% ---------------------------------------------------------
\subsection{A.8. Прокси (K, G, A\_ij): измерительные ярлыки (норматив)}

\begin{axiom}[Прокси как язык измерений, не сущности]
В целях сопоставимости расчётов, симуляций и эксперимента допускается использование прокси-наборов
\((K,\,G,\,A_{ij})\) как \emph{измерительных ярлыков} для состояния хаотической фазы в окрестности \(\Sigma\)
(т.е. для канонических понятий \(\KP\) и \(\VR\)).

\textbf{Нормативно:} прокси не являются новыми сущностями и не добавляют причинности в модель.
Причинность остаётся кинетической: соударения, КУБК (баланс числа контактов) и анизотропия приращений импульса на \(\Sigma\).
\end{axiom}

\begin{remark}[Запретная лексика для прокси]
Запрещены формулировки причинности вида: ``\(G\) действует/толкает'',
``тензор \(A_{ij}\) создаёт силу'', ``поле \(K\) ускоряет'' и т.п.
Допустимо только: ``прокси \(G\)/\(A_{ij}\) наблюдается/оценён'' и ``используется как измерительная метка''.
\end{remark}

\begin{remark}[Рекомендуемая нормировка (не обязательна)]
Допускается использовать единый вид для первых моментов распределения:
\[
K=n\langle v^2\rangle,\qquad G=n\langle v\rangle,\qquad
A_{ij}=P_{ij}-\frac{1}{3}\tr(P)\,\delta_{ij},\qquad P_{ij}=n\langle v_i v_j\rangle.
\]
\end{remark}

% ---------------------------------------------------------
\subsection{A.9. Нормы тестов шнека: ось e\_\parallel и разведение ``V=0''}

\begin{postulate}[Канонический выбор оси \(e_{\parallel}\)]
Для тестов шнека вводится ось объекта \(e_{\parallel}\) — \textbf{ось симметрии} тора/двутория,
вдоль которой определяется продольный дрейф.

\textbf{Нормативно:} \(e_{\parallel}\) должен быть однозначно воспроизводим из геометрии объекта и его паспорта,
без апелляции к внешним ``полям''.
Для тора/двутория ориентация \(e_{\parallel}\) согласуется с принятой параметризацией \(X(\varphi,\theta)\) поверхности и паспортом вихря:
касательное направление \(\partial_{\varphi}X\) соответствует циркуляции по большой окружности, а \(e_{\parallel}\) — оси симметрии выбранной системы координат.
\end{postulate}

\begin{postulate}[Разведение норм ``\(V=0\)'' и нормы по \(V_{\parallel}\)]
В тестах используются две разные нормы ``нуля'', их нельзя смешивать:
\begin{enumerate}
\item \textbf{Охранная норма (T2/T6):} запись ``\(V=0\)'' означает \textbf{нулевой вектор} скорости дрейфа,
т.е. ноль всех компонент:
\[
V=0 \iff V_x=V_y=V_z=0.
\]
\item \textbf{Норма шнека (T5--T6):} нормативной величиной является \textbf{продольная компонента}
\[
V_{\parallel}\equiv V\cdot e_{\parallel}.
\]
При прочих равных \(\chi\)-нечётной обязана быть именно \(V_{\parallel}\):
\[
V_{\parallel}(\chi)=-V_{\parallel}(-\chi).
\]
При фиксированном \(\chi\) реверс направления бега шнека обязан реверсировать знак \(V_{\parallel}\).
Для стоячего шнека нормативно \(V_{\parallel}=0\) (при отсутствии внешнего нарушителя).
\end{enumerate}
\end{postulate}

\begin{remark}
Смысл нормы: \(\chi\) должна входить в замыкания так, чтобы \(\chi\)-нечётная часть отклика давала вклад в \(V_{\parallel}\),
а не подменялась ``скрытым потоком \SE'' или асимметрией числа контактов.
\end{remark}

\section{Аксиоматическая рамка текущего шага: локальный патч как временное сужение задачи}
\label{sec:a_current_step}

\subsection{Глобальная цель и смысл временного сужения}

Цель текущего шага не состоит в замене глобальной тороидальной или двуториальной
геометрии ЭВО локальным цилиндрическим описанием.
Глобальная цель остаётся прежней:
выявить кинетический механизм дрейфа и ротации ЭВО через контактную статистику
и анизотропию приращений импульса на граничной поверхности $\Sigma$.

Квазицилиндрическое приближение вводится только как временное локальное сужение задачи.
Оно используется там, где режим тонкого тора $r \ll R$ позволяет считать короткий участок
большой окружности локально цилиндрическим.
Тем самым локальный патч $\Sigma_{\mathrm{loc}} \subset \Sigma$
служит не новой онтологией, а увеличительным языком записи,
удобным для выделения локальных механизмов на ГП.

\subsection{Что именно сохраняется неизменным}

В аксиоматическом смысле сохраняются следующие положения:

\begin{enumerate}
\item ЭВО глобально остаётся тороидальной или двуториальной структурой.
\item КУБК относится только к балансу числа контактов на $\Sigma$.
\item Дрейф и вращение допускаются только через анизотропию $\Delta p$ на $\Sigma$.
\item Шнек остаётся динамическим режимом ГП, а не отдельной сущностью и не потоком СЭ.
\item Локальный патч не заменяет глобальный интеграл по $\Sigma$, а только помогает его разложить.
\end{enumerate}

Следовательно, переход к $\Sigma_{\mathrm{loc}}$ не является сменой физической картины,
а является временным усилением разрешающей способности математического языка.

\subsection{Почему непрерывная q-схема оказалась переходной}

На предыдущем этапе непрерывный параметр $q$ использовался как язык локальной намотки.
Однако такой язык оказался недостаточным для сохранения физики шнекового процесса,
поскольку не различал явно:

\begin{itemize}
\item траектории Амеров в канале;
\item шнековую моду граничной поверхности;
\item дискретную связность шнековой организации с канальной навивкой.
\end{itemize}

Поэтому в текущем шаге физический приоритет непрерывного $q$ снимается.
Если параметр $q$ и сохраняется, то только как вспомогательный эффективный
локальный параметр записи $q_{\mathrm{eff}}$.

\subsection{Новая кандидатная рамка текущего шага}

В рабочем контуре проекта принимается кандидатная схема:

\[
m_A \in \mathbb{Z}_{>0},
\qquad
m_S \in \{1,2,4\},
\]

где $m_A$ задаёт квантовое число канальной навивки,
а $m_S$ задаёт базовую пространственную моду шнека на ГП.

Базовыми режимами считаются:
\[
S1,\quad S2,\quad S4,
\]
а все прочие режимы допускаются только как производные от них.

В локальном патч-языке должны различаться два касательных направления:
\[
h_A \quad \text{и} \quad h_S,
\]
где $h_A$ есть директор траектории Амеров,
а $h_S$ --- директор шнековой моды ГП.
Шнек не отождествляется с направлением канальной траектории.

\subsection{Интерпретационная граница}

Текущий шаг проекта следует понимать как последовательность:

\[
\text{глобальная цель}
\;\to\;
\text{тонкий тор}
\;\to\;
\Sigma_{\mathrm{loc}}
\;\to\;
(h_A,h_S)
\;\to\;
(m_A,m_S,S1/S2/S4)
\;\to\;
\text{возврат к глобальному интегралу по } \Sigma.
\]

Тем самым локальный акцент на квазицилиндрическом патче не замещает канон,
а служит промежуточным этапом прояснения механизма.
Именно в таком виде данный шаг должен отражаться в витрине и в журнале изменений.

% Совместимость: article, pdfLaTeX, T2A, babel (ru/uk)
% =========================================================

\section{Аксиоматика и определения A (расширение канона)}

\subsection*{Статус}
\textbf{Ответственный:} Сотрудник A (аксиоматика) \quad
\textbf{Статус:} рабочая вставка \quad
\textbf{Версия:} A-0.2 \quad
\textbf{Дата (Europe/Kyiv):} \today

% ---------------------------------------------------------
\subsection{A.1. Термины и фазовые состояния}

\begin{definition}[Фазовые состояния эфира]
В модели различаются два фазовых состояния движения Амеров:
\begin{itemize}
\item \textbf{Хаотическая фаза:} статистически хаотическое движение Амеров.
Вне \EVO\ хаотическая фаза называется \SE\ (Свободный эфир), а внутри \EVO\ — \NE\ (Несвободный эфир),
включая центральное ядро \NY\ и межканальные прослойки.
\item \textbf{Направленная фаза:} упорядоченное циркуляционное движение, существующее \emph{только} в каналах \EVO.
Направленное движение в каналах \textbf{не} называется \NE.
\end{itemize}
\end{definition}

\begin{axiom}[Запрет устойчивых потоков \SE]
В \SE\ не допускаются устойчивые потоки как первичная онтология.
Допускаются только распространяющиеся \VR\ как флуктуации импульса/плотности,
передаваемые цепочками парных соударений.
\end{axiom}

% ---------------------------------------------------------
\subsection{A.2. Геометрия \EVO: каналы и граничные поверхности}

\begin{definition}[Каналы и граничные поверхности (\GP)]
\EVO\ состоит из одного или нескольких коаксиальных циркуляционных каналов (направленная фаза)
и областей \NE\ (хаос) между ними и в \NY.
Каждый канал имеет две \GP: внутреннюю и внешнюю.
Внешняя \GP внешнего канала является внешней \GP всего \EVO\ (границей \EVO\ с \SE).
\end{definition}

\begin{remark}[Шнековая организация \GP]
Шнекоподобная форма \GP понимается как согласованная пространственно-временная организация контактов,
которая может быть бегущей или стоячей (см. A.6) и влияет на статистику соударений на \GP.
\end{remark}

% ---------------------------------------------------------
\subsection{A.3. Инвариант и жёсткие запреты}

\begin{axiom}[Единственный физический акт]
Единственный физический акт модели — парные упругие соударения Амеров.
Все процессы (структурирование, перенос \VR, взаимодействия, движение/вращение \EVO) реализуются только через них.
\end{axiom}

\begin{axiom}[Жёсткие запреты модели]
Запрещается вводить как первичные сущности:
\begin{itemize}
\item устойчивые потоки \SE;
\item ``давление \SE'' как первичный агент;
\item проникновение Амеров через \GP;
\item поглощение/переизлучение Амеров;
\item ``виртуальные переносчики'' как физические объекты;
\item ``поле'' как отдельную субстанцию (термин допустим только как макро-ярлык).
\end{itemize}
Генезис вихрей из \SE\ в рамках канона не рассматривается (\EVO\ считаются существующими).
\end{axiom}

% ---------------------------------------------------------
\subsection{A.4. Условие целостности хаотической фазы (запрет на обгон)}

\begin{postulate}[Целостность хаотической фазы у \GP]
В хаотической фазе (\SE\ и \NE) у \GP поддерживается режим, обеспечивающий целостность \EVO:
\[
l_{\max}\le d_A, \qquad \mean{l}=0.5\,d_A,
\]
где \(l\) — длина свободного пробега Амера, \(d_A\) — диаметр Амера.
При нарушении режима целостности вихревая структура теряет устойчивость и распадается.
\end{postulate}

\begin{remark}[Каналы]
В каналах (направленная фаза) концентрация ниже, скорости выше;
обгоны/отставания возможны внутри потока, но отсутствует движение против общего направления.
\end{remark}

% ---------------------------------------------------------
\subsection{A.5. КУБК: контактный баланс и минимальная концентрация канала}

\begin{postulate}[КУБК: баланс \emph{числа} контактов на \GP]
В квазистационарном режиме на каждом участке \GP отсутствует устойчивый перекос по \emph{числу} контактов
между амерами хаотической фазы (\SE/\NE) и амерами канала у \GP.
Эквивалентно, после осреднения по времени выполняется согласование \emph{интенсивностей контактов (касания \GP)}:
\[
\Phi^-_{\mathrm{хаос}}(x)=\Phi^+_{\mathrm{канал}}(x),\qquad x\in\Sigma,
\]
где \(\Sigma\) — соответствующая \GP, \(n\) — внешняя единичная нормаль к \(\Sigma\),
а \(\Phi\) не является ``потоком \SE\ в объёме'', а есть локальная величина \emph{частоты/интенсивности касаний поверхности}:
\[
\Phi^-_{\mathrm{хаос}}(x)=\int_{v\cdot n<0}\abs{v\cdot n}\;f_{\mathrm{хаос}}(x,v)\,d^3v,\qquad
\Phi^+_{\mathrm{канал}}(x)=\int_{v\cdot n>0}(v\cdot n)\;f_{\mathrm{канал}}(x,v)\,d^3v.
\]
Здесь \(f_{\mathrm{хаос}}(x,v)\) и \(f_{\mathrm{канал}}(x,v)\) — локальные функции распределения амеров по скоростям
в окрестности \(\Sigma\) со своей стороны поверхности; интегрирование ведётся по соответствующим полупространствам
\(v\cdot n<0\) и \(v\cdot n>0\).
\end{postulate}

\begin{postulate}[Принцип минимальной концентрации канала у \GP]
Среди всех режимов канала, совместимых с целостностью \EVO, реализуется режим с \emph{минимально возможной}
локальной концентрацией \(n_{\mathrm{канал}}(x)\) у \GP, при котором сохраняется КУБК.
\end{postulate}

\begin{remark}[Различение балансов]
КУБК фиксирует баланс \emph{числа} контактов.
Дрейф/вращение определяются не этим балансом, а анизотропией \emph{приращений импульса} в контактах на \GP.
\end{remark}

% ---------------------------------------------------------
\subsection{A.6. Динамическая природа ``шнека''}

\begin{postulate}[Шнек как колебательный режим \GP]
Равновесие на \GP поддерживается колебательным процессом границы относительно среднего положения.
Колебания \GP локально ортогональны направлению основного потока в канале.
Период/фаза колебаний связаны (или фазно согласованы) с характерным временем циркуляции в структуре канала,
в частности с временем обращения вокруг малого кольца тора и/или его гармониками.
Возможны режимы: стоячий шнек и бегущий шнек (по ходу или против хода канального движения),
в зависимости от распределения внутренних/межканальных \VR\ и состояния \KP.
\end{postulate}

% ---------------------------------------------------------
\subsection{A.7. Симметрия дрейфа, хиральность и двуторий}

\begin{postulate}[Симметрия дрейфа в изотропном фоне]
В изотропном фоне (\(\nabla\KP=0\)) и при отсутствии внешних \VR, выделяющих направление,
идеализированный симметричный ``тонкий'' тор не имеет выделенного направления для устойчивого дрейфа.
Следовательно, дрейф требует нарушения симметрии статистики \emph{приращений импульса} на \GP
(внешнего или самосогласованного внутреннего).
\end{postulate}

\begin{definition}[Хиральность (зарученность)]
Хиральность \(\chi\) — дискретный паспортный знак зарученности циркуляции канала/объекта:
\[
\chi\in\{-1,+1\}.
\]
\(\chi\) используется как метка конфигурации (паспорт вихря) и не вводится как отдельная физическая субстанция.
\end{definition}

\begin{definition}[Двуторий]
Двуторий — тороидально-спиральная форма \EVO, возникающая как самосогласованная трансформация ``тонкого'' вихря,
когда собственные \VR\ и связанная с ними структура \KP\ внутри \EVO\ приводят к устойчивому нарушению симметрии
статистики контактов/приращений импульса на \GP.
\end{definition}

\begin{postulate}[Двуторий: механизм ``вращение \(\to\) дрейф'']
В двутории внутренняя топология и шнековая организация \GP обеспечивают преобразование вращательных мод
в поступательный дрейф вдоль оси симметрии.
Механизм трактуется кинетически: интерференция собственных \VR\ формирует устойчивую анизотропию состояния хаотической фазы
(градиенты \KP) у разных участков \GP; интеграл микросмещений мест соударений на \GP даёт ненулевой дрейф и/или момент.
\end{postulate}

\begin{remark}
Формализация дрейфа/момента через поверхностные интегралы от анизотропии импульсного обмена \(j(x,t)\)
является предметом раздела B/Математика.
\end{remark}

% ---------------------------------------------------------
\subsection{A.8. Прокси (K, G, A\_ij): измерительные ярлыки (норматив)}

\begin{axiom}[Прокси как язык измерений, не сущности]
В целях сопоставимости расчётов, симуляций и эксперимента допускается использование прокси-наборов
\((K,\,G,\,A_{ij})\) как \emph{измерительных ярлыков} для состояния хаотической фазы в окрестности \(\Sigma\)
(т.е. для канонических понятий \(\KP\) и \(\VR\)).

\textbf{Нормативно:} прокси не являются новыми сущностями и не добавляют причинности в модель.
Причинность остаётся кинетической: соударения, КУБК (баланс числа контактов) и анизотропия приращений импульса на \(\Sigma\).
\end{axiom}

\begin{remark}[Запретная лексика для прокси]
Запрещены формулировки причинности вида: ``\(G\) действует/толкает'',
``тензор \(A_{ij}\) создаёт силу'', ``поле \(K\) ускоряет'' и т.п.
Допустимо только: ``прокси \(G\)/\(A_{ij}\) наблюдается/оценён'' и ``используется как измерительная метка''.
\end{remark}

\begin{remark}[Рекомендуемая нормировка (не обязательна)]
Допускается использовать единый вид для первых моментов распределения:
\[
K=n\langle v^2\rangle,\qquad G=n\langle v\rangle,\qquad
A_{ij}=P_{ij}-\frac{1}{3}\tr(P)\,\delta_{ij},\qquad P_{ij}=n\langle v_i v_j\rangle.
\]
\end{remark}

% ---------------------------------------------------------
\subsection{A.9. Нормы тестов шнека: ось e\_\parallel и разведение ``V=0''}

\begin{postulate}[Канонический выбор оси \(e_{\parallel}\)]
Для тестов шнека вводится ось объекта \(e_{\parallel}\) — \textbf{ось симметрии} тора/двутория,
вдоль которой определяется продольный дрейф.

\textbf{Нормативно:} \(e_{\parallel}\) должен быть однозначно воспроизводим из геометрии объекта и его паспорта,
без апелляции к внешним ``полям''.
Для тора/двутория ориентация \(e_{\parallel}\) согласуется с принятой параметризацией \(X(\varphi,\theta)\) поверхности и паспортом вихря:
касательное направление \(\partial_{\varphi}X\) соответствует циркуляции по большой окружности, а \(e_{\parallel}\) — оси симметрии выбранной системы координат.
\end{postulate}

\begin{postulate}[Разведение норм ``\(V=0\)'' и нормы по \(V_{\parallel}\)]
В тестах используются две разные нормы ``нуля'', их нельзя смешивать:
\begin{enumerate}
\item \textbf{Охранная норма (T2/T6):} запись ``\(V=0\)'' означает \textbf{нулевой вектор} скорости дрейфа,
т.е. ноль всех компонент:
\[
V=0 \iff V_x=V_y=V_z=0.
\]
\item \textbf{Норма шнека (T5--T6):} нормативной величиной является \textbf{продольная компонента}
\[
V_{\parallel}\equiv V\cdot e_{\parallel}.
\]
При прочих равных \(\chi\)-нечётной обязана быть именно \(V_{\parallel}\):
\[
V_{\parallel}(\chi)=-V_{\parallel}(-\chi).
\]
При фиксированном \(\chi\) реверс направления бега шнека обязан реверсировать знак \(V_{\parallel}\).
Для стоячего шнека нормативно \(V_{\parallel}=0\) (при отсутствии внешнего нарушителя).
\end{enumerate}
\end{postulate}

\begin{remark}
Смысл нормы: \(\chi\) должна входить в замыкания так, чтобы \(\chi\)-нечётная часть отклика давала вклад в \(V_{\parallel}\),
а не подменялась ``скрытым потоком \SE'' или асимметрией числа контактов.
\end{remark}

\section{Аксиоматическая рамка текущего шага: локальный патч как временное сужение задачи}
\label{sec:a_current_step}

\subsection{Глобальная цель и смысл временного сужения}

Цель текущего шага не состоит в замене глобальной тороидальной или двуториальной
геометрии ЭВО локальным цилиндрическим описанием.
Глобальная цель остаётся прежней:
выявить кинетический механизм дрейфа и ротации ЭВО через контактную статистику
и анизотропию приращений импульса на граничной поверхности $\Sigma$.

Квазицилиндрическое приближение вводится только как временное локальное сужение задачи.
Оно используется там, где режим тонкого тора $r \ll R$ позволяет считать короткий участок
большой окружности локально цилиндрическим.
Тем самым локальный патч $\Sigma_{\mathrm{loc}} \subset \Sigma$
служит не новой онтологией, а увеличительным языком записи,
удобным для выделения локальных механизмов на ГП.

\subsection{Что именно сохраняется неизменным}

В аксиоматическом смысле сохраняются следующие положения:

\begin{enumerate}
\item ЭВО глобально остаётся тороидальной или двуториальной структурой.
\item КУБК относится только к балансу числа контактов на $\Sigma$.
\item Дрейф и вращение допускаются только через анизотропию $\Delta p$ на $\Sigma$.
\item Шнек остаётся динамическим режимом ГП, а не отдельной сущностью и не потоком СЭ.
\item Локальный патч не заменяет глобальный интеграл по $\Sigma$, а только помогает его разложить.
\end{enumerate}

Следовательно, переход к $\Sigma_{\mathrm{loc}}$ не является сменой физической картины,
а является временным усилением разрешающей способности математического языка.

\subsection{Почему непрерывная q-схема оказалась переходной}

На предыдущем этапе непрерывный параметр $q$ использовался как язык локальной намотки.
Однако такой язык оказался недостаточным для сохранения физики шнекового процесса,
поскольку не различал явно:

\begin{itemize}
\item траектории Амеров в канале;
\item шнековую моду граничной поверхности;
\item дискретную связность шнековой организации с канальной навивкой.
\end{itemize}

Поэтому в текущем шаге физический приоритет непрерывного $q$ снимается.
Если параметр $q$ и сохраняется, то только как вспомогательный эффективный
локальный параметр записи $q_{\mathrm{eff}}$.

\subsection{Новая кандидатная рамка текущего шага}

В рабочем контуре проекта принимается кандидатная схема:

\[
m_A \in \mathbb{Z}_{>0},
\qquad
m_S \in \{1,2,4\},
\]

где $m_A$ задаёт квантовое число канальной навивки,
а $m_S$ задаёт базовую пространственную моду шнека на ГП.

Базовыми режимами считаются:
\[
S1,\quad S2,\quad S4,
\]
а все прочие режимы допускаются только как производные от них.

В локальном патч-языке должны различаться два касательных направления:
\[
h_A \quad \text{и} \quad h_S,
\]
где $h_A$ есть директор траектории Амеров,
а $h_S$ --- директор шнековой моды ГП.
Шнек не отождествляется с направлением канальной траектории.

\subsection{Интерпретационная граница}

Текущий шаг проекта следует понимать как последовательность:

\[
\text{глобальная цель}
\;\to\;
\text{тонкий тор}
\;\to\;
\Sigma_{\mathrm{loc}}
\;\to\;
(h_A,h_S)
\;\to\;
(m_A,m_S,S1/S2/S4)
\;\to\;
\text{возврат к глобальному интегралу по } \Sigma.
\]

Тем самым локальный акцент на квазицилиндрическом патче не замещает канон,
а служит промежуточным этапом прояснения механизма.
Именно в таком виде данный шаг должен отражаться в витрине и в журнале изменений.

% Совместимость: article, pdfLaTeX, T2A, babel (ru/uk)
% =========================================================

\section{Аксиоматика и определения A (расширение канона)}

\subsection*{Статус}
\textbf{Ответственный:} Сотрудник A (аксиоматика) \quad
\textbf{Статус:} рабочая вставка \quad
\textbf{Версия:} A-0.2 \quad
\textbf{Дата (Europe/Kyiv):} \today

% ---------------------------------------------------------
\subsection{A.1. Термины и фазовые состояния}

\begin{definition}[Фазовые состояния эфира]
В модели различаются два фазовых состояния движения Амеров:
\begin{itemize}
\item \textbf{Хаотическая фаза:} статистически хаотическое движение Амеров.
Вне \EVO\ хаотическая фаза называется \SE\ (Свободный эфир), а внутри \EVO\ — \NE\ (Несвободный эфир),
включая центральное ядро \NY\ и межканальные прослойки.
\item \textbf{Направленная фаза:} упорядоченное циркуляционное движение, существующее \emph{только} в каналах \EVO.
Направленное движение в каналах \textbf{не} называется \NE.
\end{itemize}
\end{definition}

\begin{axiom}[Запрет устойчивых потоков \SE]
В \SE\ не допускаются устойчивые потоки как первичная онтология.
Допускаются только распространяющиеся \VR\ как флуктуации импульса/плотности,
передаваемые цепочками парных соударений.
\end{axiom}

% ---------------------------------------------------------
\subsection{A.2. Геометрия \EVO: каналы и граничные поверхности}

\begin{definition}[Каналы и граничные поверхности (\GP)]
\EVO\ состоит из одного или нескольких коаксиальных циркуляционных каналов (направленная фаза)
и областей \NE\ (хаос) между ними и в \NY.
Каждый канал имеет две \GP: внутреннюю и внешнюю.
Внешняя \GP внешнего канала является внешней \GP всего \EVO\ (границей \EVO\ с \SE).
\end{definition}

\begin{remark}[Шнековая организация \GP]
Шнекоподобная форма \GP понимается как согласованная пространственно-временная организация контактов,
которая может быть бегущей или стоячей (см. A.6) и влияет на статистику соударений на \GP.
\end{remark}

% ---------------------------------------------------------
\subsection{A.3. Инвариант и жёсткие запреты}

\begin{axiom}[Единственный физический акт]
Единственный физический акт модели — парные упругие соударения Амеров.
Все процессы (структурирование, перенос \VR, взаимодействия, движение/вращение \EVO) реализуются только через них.
\end{axiom}

\begin{axiom}[Жёсткие запреты модели]
Запрещается вводить как первичные сущности:
\begin{itemize}
\item устойчивые потоки \SE;
\item ``давление \SE'' как первичный агент;
\item проникновение Амеров через \GP;
\item поглощение/переизлучение Амеров;
\item ``виртуальные переносчики'' как физические объекты;
\item ``поле'' как отдельную субстанцию (термин допустим только как макро-ярлык).
\end{itemize}
Генезис вихрей из \SE\ в рамках канона не рассматривается (\EVO\ считаются существующими).
\end{axiom}

% ---------------------------------------------------------
\subsection{A.4. Условие целостности хаотической фазы (запрет на обгон)}

\begin{postulate}[Целостность хаотической фазы у \GP]
В хаотической фазе (\SE\ и \NE) у \GP поддерживается режим, обеспечивающий целостность \EVO:
\[
l_{\max}\le d_A, \qquad \mean{l}=0.5\,d_A,
\]
где \(l\) — длина свободного пробега Амера, \(d_A\) — диаметр Амера.
При нарушении режима целостности вихревая структура теряет устойчивость и распадается.
\end{postulate}

\begin{remark}[Каналы]
В каналах (направленная фаза) концентрация ниже, скорости выше;
обгоны/отставания возможны внутри потока, но отсутствует движение против общего направления.
\end{remark}

% ---------------------------------------------------------
\subsection{A.5. КУБК: контактный баланс и минимальная концентрация канала}

\begin{postulate}[КУБК: баланс \emph{числа} контактов на \GP]
В квазистационарном режиме на каждом участке \GP отсутствует устойчивый перекос по \emph{числу} контактов
между амерами хаотической фазы (\SE/\NE) и амерами канала у \GP.
Эквивалентно, после осреднения по времени выполняется согласование \emph{интенсивностей контактов (касания \GP)}:
\[
\Phi^-_{\mathrm{хаос}}(x)=\Phi^+_{\mathrm{канал}}(x),\qquad x\in\Sigma,
\]
где \(\Sigma\) — соответствующая \GP, \(n\) — внешняя единичная нормаль к \(\Sigma\),
а \(\Phi\) не является ``потоком \SE\ в объёме'', а есть локальная величина \emph{частоты/интенсивности касаний поверхности}:
\[
\Phi^-_{\mathrm{хаос}}(x)=\int_{v\cdot n<0}\abs{v\cdot n}\;f_{\mathrm{хаос}}(x,v)\,d^3v,\qquad
\Phi^+_{\mathrm{канал}}(x)=\int_{v\cdot n>0}(v\cdot n)\;f_{\mathrm{канал}}(x,v)\,d^3v.
\]
Здесь \(f_{\mathrm{хаос}}(x,v)\) и \(f_{\mathrm{канал}}(x,v)\) — локальные функции распределения амеров по скоростям
в окрестности \(\Sigma\) со своей стороны поверхности; интегрирование ведётся по соответствующим полупространствам
\(v\cdot n<0\) и \(v\cdot n>0\).
\end{postulate}

\begin{postulate}[Принцип минимальной концентрации канала у \GP]
Среди всех режимов канала, совместимых с целостностью \EVO, реализуется режим с \emph{минимально возможной}
локальной концентрацией \(n_{\mathrm{канал}}(x)\) у \GP, при котором сохраняется КУБК.
\end{postulate}

\begin{remark}[Различение балансов]
КУБК фиксирует баланс \emph{числа} контактов.
Дрейф/вращение определяются не этим балансом, а анизотропией \emph{приращений импульса} в контактах на \GP.
\end{remark}

% ---------------------------------------------------------
\subsection{A.6. Динамическая природа ``шнека''}

\begin{postulate}[Шнек как колебательный режим \GP]
Равновесие на \GP поддерживается колебательным процессом границы относительно среднего положения.
Колебания \GP локально ортогональны направлению основного потока в канале.
Период/фаза колебаний связаны (или фазно согласованы) с характерным временем циркуляции в структуре канала,
в частности с временем обращения вокруг малого кольца тора и/или его гармониками.
Возможны режимы: стоячий шнек и бегущий шнек (по ходу или против хода канального движения),
в зависимости от распределения внутренних/межканальных \VR\ и состояния \KP.
\end{postulate}

% ---------------------------------------------------------
\subsection{A.7. Симметрия дрейфа, хиральность и двуторий}

\begin{postulate}[Симметрия дрейфа в изотропном фоне]
В изотропном фоне (\(\nabla\KP=0\)) и при отсутствии внешних \VR, выделяющих направление,
идеализированный симметричный ``тонкий'' тор не имеет выделенного направления для устойчивого дрейфа.
Следовательно, дрейф требует нарушения симметрии статистики \emph{приращений импульса} на \GP
(внешнего или самосогласованного внутреннего).
\end{postulate}

\begin{definition}[Хиральность (зарученность)]
Хиральность \(\chi\) — дискретный паспортный знак зарученности циркуляции канала/объекта:
\[
\chi\in\{-1,+1\}.
\]
\(\chi\) используется как метка конфигурации (паспорт вихря) и не вводится как отдельная физическая субстанция.
\end{definition}

\begin{definition}[Двуторий]
Двуторий — тороидально-спиральная форма \EVO, возникающая как самосогласованная трансформация ``тонкого'' вихря,
когда собственные \VR\ и связанная с ними структура \KP\ внутри \EVO\ приводят к устойчивому нарушению симметрии
статистики контактов/приращений импульса на \GP.
\end{definition}

\begin{postulate}[Двуторий: механизм ``вращение \(\to\) дрейф'']
В двутории внутренняя топология и шнековая организация \GP обеспечивают преобразование вращательных мод
в поступательный дрейф вдоль оси симметрии.
Механизм трактуется кинетически: интерференция собственных \VR\ формирует устойчивую анизотропию состояния хаотической фазы
(градиенты \KP) у разных участков \GP; интеграл микросмещений мест соударений на \GP даёт ненулевой дрейф и/или момент.
\end{postulate}

\begin{remark}
Формализация дрейфа/момента через поверхностные интегралы от анизотропии импульсного обмена \(j(x,t)\)
является предметом раздела B/Математика.
\end{remark}

% ---------------------------------------------------------
\subsection{A.8. Прокси (K, G, A\_ij): измерительные ярлыки (норматив)}

\begin{axiom}[Прокси как язык измерений, не сущности]
В целях сопоставимости расчётов, симуляций и эксперимента допускается использование прокси-наборов
\((K,\,G,\,A_{ij})\) как \emph{измерительных ярлыков} для состояния хаотической фазы в окрестности \(\Sigma\)
(т.е. для канонических понятий \(\KP\) и \(\VR\)).

\textbf{Нормативно:} прокси не являются новыми сущностями и не добавляют причинности в модель.
Причинность остаётся кинетической: соударения, КУБК (баланс числа контактов) и анизотропия приращений импульса на \(\Sigma\).
\end{axiom}

\begin{remark}[Запретная лексика для прокси]
Запрещены формулировки причинности вида: ``\(G\) действует/толкает'',
``тензор \(A_{ij}\) создаёт силу'', ``поле \(K\) ускоряет'' и т.п.
Допустимо только: ``прокси \(G\)/\(A_{ij}\) наблюдается/оценён'' и ``используется как измерительная метка''.
\end{remark}

\begin{remark}[Рекомендуемая нормировка (не обязательна)]
Допускается использовать единый вид для первых моментов распределения:
\[
K=n\langle v^2\rangle,\qquad G=n\langle v\rangle,\qquad
A_{ij}=P_{ij}-\frac{1}{3}\tr(P)\,\delta_{ij},\qquad P_{ij}=n\langle v_i v_j\rangle.
\]
\end{remark}

% ---------------------------------------------------------
\subsection{A.9. Нормы тестов шнека: ось e\_\parallel и разведение ``V=0''}

\begin{postulate}[Канонический выбор оси \(e_{\parallel}\)]
Для тестов шнека вводится ось объекта \(e_{\parallel}\) — \textbf{ось симметрии} тора/двутория,
вдоль которой определяется продольный дрейф.

\textbf{Нормативно:} \(e_{\parallel}\) должен быть однозначно воспроизводим из геометрии объекта и его паспорта,
без апелляции к внешним ``полям''.
Для тора/двутория ориентация \(e_{\parallel}\) согласуется с принятой параметризацией \(X(\varphi,\theta)\) поверхности и паспортом вихря:
касательное направление \(\partial_{\varphi}X\) соответствует циркуляции по большой окружности, а \(e_{\parallel}\) — оси симметрии выбранной системы координат.
\end{postulate}

\begin{postulate}[Разведение норм ``\(V=0\)'' и нормы по \(V_{\parallel}\)]
В тестах используются две разные нормы ``нуля'', их нельзя смешивать:
\begin{enumerate}
\item \textbf{Охранная норма (T2/T6):} запись ``\(V=0\)'' означает \textbf{нулевой вектор} скорости дрейфа,
т.е. ноль всех компонент:
\[
V=0 \iff V_x=V_y=V_z=0.
\]
\item \textbf{Норма шнека (T5--T6):} нормативной величиной является \textbf{продольная компонента}
\[
V_{\parallel}\equiv V\cdot e_{\parallel}.
\]
При прочих равных \(\chi\)-нечётной обязана быть именно \(V_{\parallel}\):
\[
V_{\parallel}(\chi)=-V_{\parallel}(-\chi).
\]
При фиксированном \(\chi\) реверс направления бега шнека обязан реверсировать знак \(V_{\parallel}\).
Для стоячего шнека нормативно \(V_{\parallel}=0\) (при отсутствии внешнего нарушителя).
\end{enumerate}
\end{postulate}

\begin{remark}
Смысл нормы: \(\chi\) должна входить в замыкания так, чтобы \(\chi\)-нечётная часть отклика давала вклад в \(V_{\parallel}\),
а не подменялась ``скрытым потоком \SE'' или асимметрией числа контактов.
\end{remark}

\section{Аксиоматическая рамка текущего шага: локальный патч как временное сужение задачи}
\label{sec:a_current_step}

\subsection{Глобальная цель и смысл временного сужения}

Цель текущего шага не состоит в замене глобальной тороидальной или двуториальной
геометрии ЭВО локальным цилиндрическим описанием.
Глобальная цель остаётся прежней:
выявить кинетический механизм дрейфа и ротации ЭВО через контактную статистику
и анизотропию приращений импульса на граничной поверхности $\Sigma$.

Квазицилиндрическое приближение вводится только как временное локальное сужение задачи.
Оно используется там, где режим тонкого тора $r \ll R$ позволяет считать короткий участок
большой окружности локально цилиндрическим.
Тем самым локальный патч $\Sigma_{\mathrm{loc}} \subset \Sigma$
служит не новой онтологией, а увеличительным языком записи,
удобным для выделения локальных механизмов на ГП.

\subsection{Что именно сохраняется неизменным}

В аксиоматическом смысле сохраняются следующие положения:

\begin{enumerate}
\item ЭВО глобально остаётся тороидальной или двуториальной структурой.
\item КУБК относится только к балансу числа контактов на $\Sigma$.
\item Дрейф и вращение допускаются только через анизотропию $\Delta p$ на $\Sigma$.
\item Шнек остаётся динамическим режимом ГП, а не отдельной сущностью и не потоком СЭ.
\item Локальный патч не заменяет глобальный интеграл по $\Sigma$, а только помогает его разложить.
\end{enumerate}

Следовательно, переход к $\Sigma_{\mathrm{loc}}$ не является сменой физической картины,
а является временным усилением разрешающей способности математического языка.

\subsection{Почему непрерывная q-схема оказалась переходной}

На предыдущем этапе непрерывный параметр $q$ использовался как язык локальной намотки.
Однако такой язык оказался недостаточным для сохранения физики шнекового процесса,
поскольку не различал явно:

\begin{itemize}
\item траектории Амеров в канале;
\item шнековую моду граничной поверхности;
\item дискретную связность шнековой организации с канальной навивкой.
\end{itemize}

Поэтому в текущем шаге физический приоритет непрерывного $q$ снимается.
Если параметр $q$ и сохраняется, то только как вспомогательный эффективный
локальный параметр записи $q_{\mathrm{eff}}$.

\subsection{Новая кандидатная рамка текущего шага}

В рабочем контуре проекта принимается кандидатная схема:

\[
m_A \in \mathbb{Z}_{>0},
\qquad
m_S \in \{1,2,4\},
\]

где $m_A$ задаёт квантовое число канальной навивки,
а $m_S$ задаёт базовую пространственную моду шнека на ГП.

Базовыми режимами считаются:
\[
S1,\quad S2,\quad S4,
\]
а все прочие режимы допускаются только как производные от них.

В локальном патч-языке должны различаться два касательных направления:
\[
h_A \quad \text{и} \quad h_S,
\]
где $h_A$ есть директор траектории Амеров,
а $h_S$ --- директор шнековой моды ГП.
Шнек не отождествляется с направлением канальной траектории.

\subsection{Интерпретационная граница}

Текущий шаг проекта следует понимать как последовательность:

\[
\text{глобальная цель}
\;\to\;
\text{тонкий тор}
\;\to\;
\Sigma_{\mathrm{loc}}
\;\to\;
(h_A,h_S)
\;\to\;
(m_A,m_S,S1/S2/S4)
\;\to\;
\text{возврат к глобальному интегралу по } \Sigma.
\]

Тем самым локальный акцент на квазицилиндрическом патче не замещает канон,
а служит промежуточным этапом прояснения механизма.
Именно в таком виде данный шаг должен отражаться в витрине и в журнале изменений.
